\chapter{Le feed handler}\label{ch:ll:the-feed-handler}

À l'ouverture, le débit de messages bondit~: sur l'instrument actif du simulateur de ce chapitre, on passe d'environ 570 messages par seconde à
midi à plus de 9\,000 dans les premières secondes~; sur le flux consolidé des options américaines, OPRA a publié un pic de 63.9
millions de messages en une seconde en juillet 2026, et 350\,000 dans sa milliseconde la plus chargée. Un \omterm{def:ll:the-feed-handler:handler}{feed handler} qui prend du retard laisse son
tampon de socket se remplir~; le noyau jette ce qui ne tient pas~; la ligne A perd un paquet que la ligne B a perdu aussi~; et le carnet de la société
reste périmé tant que le \omterm{def:ll:the-feed-handler:handler}{feed handler} ne s'est pas rétabli, par retransmission si le trou est petit et récent, à partir du prochain instantané sinon.
Ce chapitre construit le premier composant du système de trading de la partie~IV, le \omterm{def:ll:the-feed-handler:handler}{feed handler}~: il décode et arbitre les deux
lignes du chapitre~16, détecte les trous et s'en remet, normalise chaque message en enregistrement d'événement de la société et le publie
sur les anneaux du chapitre~12, et indique à chaque instant si le carnet qu'il alimente peut être tenu pour fiable.

\section{Décoder et arbitrer}

\begin{definition}[Feed handler]\label{def:ll:the-feed-handler:handler}
Un \emph{feed handler}\index{feed handler} est le processus qui reçoit le flux de données de marché d'une plateforme (livre~1,
chapitre~4) et le transforme en événements internes de la société~: il lit les lignes redondantes, les arbitre en un flux unique
ordonné, détecte et récupère les messages perdus, normalise chaque message dans l'enregistrement propre de la société (livre~1, chapitre~28), et
publie les enregistrements à destination des composants qui construisent les carnets et prennent les décisions.
\end{definition}

\begin{omfigure}
\begin{tikzpicture}[font=\footnotesize,box/.style={draw,rounded corners=2pt,minimum height=0.7cm,align=center,inner sep=3pt},
  >=Stealth]
\node[box,fill=gray!10] (a) at (0,0.6) {ligne A};
\node[box,fill=gray!10] (b) at (0,-0.6) {ligne B};
\node[box,fill=omDef!20,minimum width=2.1cm] (arb) at (2.6,0) {arbitrer par\\ numéro de séquence};
\node[box,fill=omThm!15,minimum width=2.2cm] (gap) at (5.4,0) {trou~? attendre\\ l'autre ligne\\ \qty{0.5}{\milli\second}};
\node[box,fill=omDef!20,minimum width=2.1cm] (norm) at (8.3,0) {normaliser\\ (événement\\ de 48~octets)};
\node[box,fill=omMeth!20] (ring) at (10.9,0) {firm.ring\\ (chapitre~12)};
\node[box,fill=omThm!8] (retx) at (5.4,-2.0) {serveur de\\ retransmission};
\node[box,fill=omThm!8] (snap) at (8.3,-2.0) {canal d'instantanés\\ (snapshot)};
\node[box,fill=gray!8,align=left] (flag) at (5.4,1.7) {indicateur de carnet périmé, compteurs,\\ intervalles de péremption};
\draw[->] (a) -- (arb);
\draw[->] (b) -- (arb);
\draw[->] (arb) -- (gap);
\draw[->] (gap) -- (norm);
\draw[->] (norm) -- (ring);
\draw[->,omThm] (gap) -- node[left,font=\scriptsize]{demande} (retx);
\draw[->,omThm,dashed] (retx.east) -- node[below,font=\scriptsize]{} ++(0.4,0) |- ([yshift=-0.2cm]gap.east);
\draw[->,omThm] (snap) -- node[right,font=\scriptsize]{si le serveur ne peut pas} (norm);
\draw[gray,dashed] (gap) -- (flag);
\end{tikzpicture}

\omcaption{Le \omterm{def:ll:the-feed-handler:handler}{feed handler} de la construction. Les paquets des deux lignes sont fusionnés par numéro de séquence~; un trou retient les messages suivants et
attend brièvement l'autre ligne, puis interroge le serveur de retransmission, puis se replie sur le prochain instantané~; chaque message
sort sous la forme d'un événement de taille fixe sur l'anneau de la société. Tant qu'un trou est ouvert, l'\omterm{def:ll:the-feed-handler:stale}{indicateur de carnet périmé} est levé.}\label{fig:ll:the-feed-handler:design}
\end{omfigure}

Le \omterm{def:ll:the-feed-handler:handler}{feed handler} de la construction est écrit d'abord en Python, comme référence, puis en C++20 et en Rust, et les trois publient exactement
les mêmes événements (comparés par un condensat de 64~bits de leurs enregistrements de 48~octets) sur les mêmes entrées. Son temps est l'heure d'arrivée enregistrée
avec chaque paquet, de sorte qu'une journée enregistrée se rejoue à l'identique~: mêmes entrées, mêmes événements, mêmes intervalles de péremption
et mêmes compteurs. L'arbitrage est celui du chapitre~16~; ce que le \omterm{def:ll:the-feed-handler:handler}{feed handler} ajoute, c'est la conduite à tenir pendant qu'un numéro de séquence manque.

\section{Détecter les trous et s'en remettre}

Un numéro de séquence manquant n'est en général pas perdu, seulement en retard~: il est en route sur l'autre ligne, ou il a été réordonné par un commutateur.
Le \omterm{def:ll:the-feed-handler:handler}{feed handler} retient donc les messages qui arrivent après lui et attend un court instant, ici \qty{0.5}{\milli\second}, les
messages manquants. S'ils arrivent, les messages retenus sont libérés dans l'ordre. Sinon, le \omterm{def:ll:the-feed-handler:handler}{feed handler} demande au serveur de retransmission
la plage concernée (chapitre~16~; le simulateur répond à hauteur de 1\,000 messages par requête, à partir de ses 100\,000 derniers), et applique la réponse
à son arrivée, un aller-retour plus tard. Si le serveur ne peut pas aider, parce que le trou est plus ancien que sa fenêtre ou que le serveur est
injoignable, le \omterm{def:ll:the-feed-handler:handler}{feed handler} attend le prochain instantané de l'instrument, le charge, écarte les messages retenus qu'il reflète déjà
et reprend à partir du numéro de séquence que cet instantané donne, procédure que prennent en charge les messages d'instantané du CME avec leur dernier
numéro de séquence incrémental traité.

\omcode{code/firm/feedhandler/cpp/firm_feedhandler.hpp}{213}{237}{Chaque paquet~: le rejeter si tous ses messages sont anciens, délivrer le
message attendu suivant, retenir les autres~; un message retenu ouvre un trou et lance sa minuterie.}\label{lst:ll:the-feed-handler:incremental}

Le chemin de l'instantané est le plus délicat~: l'instantané doit être au moins aussi récent que le carnet, les messages retenus qu'il reflète déjà
doivent disparaître, et un trou qui le suit doit être de nouveau détecté.

\omcode{code/firm/feedhandler/firm_feedhandler.py}{190}{204}{Chargement d'un instantané dans la référence Python.}\label{lst:ll:the-feed-handler:snapshot}

\begin{definition}[Indicateur de carnet périmé, latence de reprise]\label{def:ll:the-feed-handler:stale}
L'\emph{indicateur de carnet périmé}\index{indicateur de carnet périmé} est la déclaration publique du \omterm{def:ll:the-feed-handler:handler}{feed handler} selon laquelle le carnet construit à partir de ses événements peut
différer de celui de la plateforme~; il est levé lorsqu'un trou est détecté et baissé lorsqu'il est comblé. La
\emph{latence de reprise}\index{latence de reprise} d'un trou est le temps écoulé entre sa détection et le moment où le carnet redevient exact~: le retard de l'autre ligne,
un aller-retour de retransmission, ou l'attente et le chargement d'un instantané.
\end{definition}

La \Cref{fig:ll:the-feed-handler:minute} montre le \omterm{def:ll:the-feed-handler:handler}{feed handler} à l'œuvre sur une minute simulée depuis l'ouverture, les deux lignes étant dégradées~: 0.2~\% des paquets
perdus indépendamment sur chaque ligne, des rafales de pertes, une \omterm{def:ll:where-latency-comes-from:tail}{gigue} (jitter) qui réordonne les paquets, et une coupure sur chaque ligne
se chevauchant pendant 15 millisecondes à 40 secondes. Sur 4\,232 trous, 4\,224 sont comblés par l'autre ligne
dans la demi-milliseconde~; le serveur de retransmission referme les huit autres en \qty{2.7}{\milli\second}. Sans le serveur, sept trous
exigent un instantané, et avec un instantané toutes les secondes le carnet reste périmé de 0.3 à 1 seconde à chaque fois, soit cinq
secondes de la minute au total. La péremption n'est alors pas une propriété du réseau, mais celle du chemin de reprise.

\begin{omfigure}
\begin{tikzpicture}
\begin{groupplot}[group style={group size=1 by 2,vertical sep=0.6cm,xlabels at=edge bottom},width=12cm,xmin=0,xmax=60,
  tick label style={font=\small},label style={font=\small},grid=major,grid style={gray!25},table/col sep=comma]
\nextgroupplot[height=3.6cm,ymin=0,ymax=1000,ylabel={msg / 100 ms},xticklabels={}]
\addplot[omDef,const plot,fill=omDef!25,draw=omDef] table[x=t_s,y=messages]{figdata/low-latency/18-the-feed-handler/rate.csv} \closedcycle;
\nextgroupplot[height=3.2cm,ymin=0.4,ymax=2.6,ytick={1,2},yticklabels={{instantané seul},{avec retrans.}},
  xlabel={secondes après l'ouverture},unbounded coords=jump]
\addplot[omMeth,line width=6pt,mark=|,mark size=4pt,mark options={line width=1.2pt}] table[x=t_s,y=retx]{figdata/low-latency/18-the-feed-handler/stale_bars.csv};
\addplot[omThm,line width=6pt] table[x=t_s,y=snapshot]{figdata/low-latency/18-the-feed-handler/stale_bars.csv};
\end{groupplot}
\end{tikzpicture}

\omcaption{Une minute simulée depuis l'ouverture. En haut~: messages publiés par tranche de 100~ms, seize fois plus nombreux à
l'ouverture qu'à midi. En bas~: les trous que l'autre ligne n'a pas pu combler, avec le serveur de retransmission (chacun refermé en moins de
\qty{3}{\milli\second}~: une graduation) et avec la seule reprise par instantané (un instantané par seconde~: une barre entre la détection et la reprise)~; les
4\,224 trous comblés par l'autre ligne sont trop courts pour être dessinés. Données~: \texttt{fig\_stale.py} (déterministe~: le simulateur et le \omterm{def:ll:the-feed-handler:handler}{feed handler}
le sont).}\label{fig:ll:the-feed-handler:minute}
\end{omfigure}

\section{Normaliser et publier}

Chaque message sort du \omterm{def:ll:the-feed-handler:handler}{feed handler} sous la forme d'un événement de 48~octets~: son genre, son côté, son instrument, son numéro de séquence, son horodatage, sa référence d'ordre,
une seconde référence (la nouvelle référence d'un remplacement, le numéro de correspondance d'une exécution, le numéro de séquence d'un instantané), son prix et sa quantité.
Un enregistrement normalisé présente trois qualités en aval. Sa taille est fixe, ce qui le fait tenir dans un emplacement d'anneau
(chapitre~12) et permet à un consommateur de le lire sans le décoder. Il est le même pour chaque plateforme que la société lit, de sorte que le constructeur de carnet
(chapitre~19) et les stratégies (chapitre~20) ne soient écrits qu'une fois. Et il porte le numéro de séquence de la plateforme, si bien qu'un consommateur peut
savoir exactement quels événements il a vus, et les événements d'un instantané peuvent indiquer quel numéro de séquence ils reflètent.
Le test de la construction publie une exécution complète sur un anneau SPSC d'emplacements de 64~octets et la relit, condensat par condensat.

Sur le portable, le gestionnaire C++ prend environ \qty{39}{\nano\second} par paquet à la médiane, \qty{160}{\nano\second} au p99
et \qty{290}{\nano\second} au p99.9, mesurés sur la minute de la \cref{fig:ll:the-feed-handler:minute} avec l'horloge lue
autour de chaque entrée~; les rares cas lents sont les trous, dont les messages retenus passent par une table ordonnée, et les réponses de retransmission.
À 9\,000 messages par seconde, le gestionnaire est occupé environ 0.04~\% du temps~: la rafale de cet instrument simulé n'est rien
pour lui. Les vraies rafales sont celles de l'\cref{dat:ll:the-feed-handler:opra}.

\begin{dated}{2026-09}{Dans quelle mesure un flux est-il sujet aux rafales~?}\label{dat:ll:the-feed-handler:opra}
Selon les métriques opérationnelles publiées par OPRA pour le flux consolidé d'options des États-Unis (rapport d'août 2026), on relève pour juillet 2026 un pic de
63.9 millions de messages en une seconde, 7.9 millions en 100 millisecondes, 0.89 million en 10 millisecondes et 0.35 million en
une milliseconde~: la milliseconde la plus chargée atteint 350 millions par seconde, soit cinq fois et demie le débit de la seconde la plus chargée. L'avis de capacité d'OPRA
explique pourquoi les abonnés prévoient des pics sur 10 millisecondes («~reflètent l'utilisation du système lors de rafales de trafic~»), et
que la prise en compte des deux flux redondants double la bande passante.
\end{dated}

\section{Rafales, tampons et consommateurs lents}

\begin{definition}[Rafale de messages, débordement du tampon de réception]\label{def:ll:the-feed-handler:burst}
Une \emph{rafale de messages}\index{rafale de messages} est une période pendant laquelle le débit d'un flux dépasse le débit que son récepteur peut soutenir. Un
\emph{débordement du tampon de réception}\index{débordement du tampon de réception} est la perte des datagrammes qui arrivent alors que le tampon de réception de la
socket est plein~: le noyau les jette sans bruit, et le gestionnaire ne voit qu'un trou.
\end{definition}

Pendant une rafale, l'écart entre le débit d'arrivée et le débit de service du gestionnaire s'accumule dans le tampon de réception de la socket,
la file M/M/1 du livre~4, chapitre~8, avec une limite stricte~: ce qui ne tient pas est jeté. La taille du tampon
se règle avec \texttt{SO\_RCVBUF}, et deux détails de Linux déterminent ce qu'il contient. Le noyau «~double cette valeur (pour
laisser de la place aux frais de gestion)~» et renvoie la valeur doublée (socket(7)), plafonnée par \texttt{net.core.rmem\_max},
soit 4~MiB sur le portable, de sorte que demander 4~MiB en obtient 8. Et le tampon est décompté pour l'empreinte mémoire totale de chaque datagramme,
et non pour sa charge utile. La \Cref{fig:ll:the-feed-handler:rcvbuf} mesure ce qu'un lecteur à l'arrêt conserve~: 50\,000 datagrammes envoyés sur
la boucle locale pendant que personne ne lit, puis la socket vidée. Chaque datagramme de 64~octets coûte 960~octets de tampon et chaque datagramme de 1\,000~octets
en coûte 2\,305~: le plus grand tampon que le portable accorde retient 8\,738 petits datagrammes, soit un quarantième des messages de
la milliseconde la plus chargée d'OPRA.

\begin{omfigure}
\begin{tikzpicture}
\begin{axis}[width=11.5cm,height=5.6cm,xmode=log,ymode=log,log basis x=2,xmin=20000,xmax=12000000,ymin=8,ymax=20000,
  xlabel={tampon accordé par le noyau (octets, échelle log~; double de la taille demandée)},ylabel={datagrammes conservés},
  tick label style={font=\small},label style={font=\small},grid=major,grid style={gray!25},
  xtick={32768,131072,524288,2097152,8388608},xticklabels={32~KiB,128~KiB,512~KiB,2~MiB,8~MiB},
  legend image code/.code={\draw[#1] (0cm,0cm) -- (0.6cm,0cm);},
  legend columns=2,legend style={font=\footnotesize,at={(0.5,-0.3)},anchor=north,draw=none,fill=none,
  /tikz/every even column/.append style={column sep=8pt}},table/col sep=comma]
\addplot[omDef,very thick,mark=*,mark size=1.4pt] table[x=granted,y=held,restrict expr to domain={\thisrow{payload}}{0:100}]{figdata/low-latency/18-the-feed-handler/measured_rcvbuf.csv};
\addplot[omThm,very thick,mark=square*,mark size=1.4pt] table[x=granted,y=held,restrict expr to domain={\thisrow{payload}}{500:2000}]{figdata/low-latency/18-the-feed-handler/measured_rcvbuf.csv};
\legend{datagrammes de 64~octets,datagrammes de 1\,000~octets}
\end{axis}
\end{tikzpicture}

\omcaption{Nombre de datagrammes qu'une socket UDP conserve pour un lecteur qui ne lit pas, selon la taille du tampon de réception~: 50\,000 datagrammes
envoyés sur la boucle locale pendant l'arrêt. Le noyau accorde le double de la taille demandée et décompte à chaque datagramme son
empreinte mémoire complète. Mesuré sur un portable (Intel Core Ultra 7 155H) sous WSL2 (noyau 6.18)~; \texttt{rmem\_max} 4~MiB. Données~:
\texttt{bench\_feed.py}.}\label{fig:ll:the-feed-handler:rcvbuf}
\end{omfigure}

Le gestionnaire est à son tour un producteur. Ses consommateurs lisent l'anneau à leur propre rythme, et un consommateur plus lent que le flux (un
\omterm{def:ll:queues-ring-buffers-and-shared-memory:bp}{consommateur lent}, chapitre~12) ne doit pas ralentir le gestionnaire~: l'anneau de diffusion n'attend jamais, et un consommateur dépassé
se resynchronise. Les deux défenses contre les rafales sont donc les mêmes aux deux extrémités du gestionnaire~: un tampon suffisant pour la
rafale que les données laissent prévoir, et un moyen de reprendre lorsque même cela ne suffit pas.

\section{Surveiller un flux}

Un gestionnaire qui se trompe en silence est pire qu'un gestionnaire qui s'arrête. La construction compte les paquets, les messages, les doublons, les trous, les trous
comblés par l'autre ligne, les demandes de retransmission et les reprises par instantané, et conserve chaque intervalle de péremption avec son début
et sa fin. Quatre de ces chiffres méritent des alertes~: le taux de trous sur chaque ligne (une ligne qui se dégrade, tandis que l'autre la masque), les trous
comblés par aucune des deux lignes (une cause commune, chapitre~16), toute reprise par instantané (un carnet faux pendant jusqu'à un
intervalle d'instantané), et la marque de pic de la file de réception (une rafale plus proche du tampon que prévu). L'\omterm{def:ll:the-feed-handler:stale}{indicateur de carnet périmé} lui-même est transmis à
chaque consommateur avec chaque événement, de sorte qu'une stratégie ne cote jamais sur un carnet que son gestionnaire sait faux.

\section{Tutoriel~: le gestionnaire sur des lignes dégradées}\label{tut:ll:the-feed-handler:tut}

\begin{tutorial}
\textbf{Objectif.} Exécuter le gestionnaire dans trois langages sur les mêmes lignes dégradées, reprendre par retransmission puis par instantané, et
vérifier le résultat par rapport à une exécution propre. \textbf{État final~:} les \cref{fig:ll:the-feed-handler:minute,fig:ll:the-feed-handler:rcvbuf},
les quantiles de temps de traitement et des tests au vert en Python, C++ et Rust.
\begin{tutsteps}
\item \textbf{Construire le jeu de données.} \texttt{make\_feed\_fixtures.py} exécute le simulateur du livre~10 pendant trois secondes avec les deux lignes
dégradées et écrit les lignes, le flux sans perte (que le serveur de retransmission conserve) et le \omterm{def:ll:protocols-ii-binary-exchange-protocols:channels}{canal d'instantanés}.
\item \textbf{Trois exécutions, trois langages.} Exécution propre, exécution dégradée avec le serveur, exécution dégradée sans lui~: la référence Python écrit
les compteurs, condensats et péremptions attendus, et les tests C++ et Rust les reproduisent à l'identique.
\item \textbf{Vérifier ce que la reprise signifie.} Avec la retransmission, l'exécution dégradée publie exactement les événements de l'exécution propre. Avec
un instantané, elle publie l'instantané à la place des messages perdus, puis exactement les événements de l'exécution propre qui suivent, et le carnet
à la fin de l'exécution est le même.
\item \textbf{Mesurer} avec \texttt{python bench\_feed.py}~: la minute de la \cref{fig:ll:the-feed-handler:minute} traitée par le
gestionnaire C++ avec l'horloge lue autour de chaque entrée (le script vérifie qu'il publie ce que publie la référence), et
l'expérience du lecteur à l'arrêt de la \cref{fig:ll:the-feed-handler:rcvbuf}.
\end{tutsteps}
\textbf{Pour aller plus loin.} Réduire l'attente de l'autre ligne à \qty{50}{\micro\second} et compter les demandes de retransmission
qui en découlent~; allonger l'intervalle d'instantané à cinq secondes et relire la péremption.
\end{tutorial}

\section{Construction~: le feed handler}\label{bld:ll:the-feed-handler:feedhandler}

\begin{build}
\bfield{Rôle} La première étape du système de trading~: des lignes de la plateforme aux événements de la société. Le constructeur de carnet du
chapitre~19 consomme ses événements, le moteur de stratégies du chapitre~20 son \omterm{def:ll:the-feed-handler:stale}{indicateur de carnet périmé}, et la surveillance du chapitre~24 ses
compteurs.
\bfield{Interface} Référence Python \texttt{firm\_feedhandler}~: \texttt{Handler(retx, gap\_timeout\_ns, rtt\_ns).run(events)},
\texttt{.events}, \texttt{.hash}, \texttt{.stale}, \texttt{.counters}~; \texttt{RetxServer}, \texttt{merge}, \texttt{recorded},
\texttt{normalise}, \texttt{book}. C++20 \texttt{firm::feed2::Handler(retx, timeout, rtt)} avec \texttt{run(a, b, snapshots)},
\texttt{on\_event}, \texttt{hash}, \texttt{stale}, compteurs et points d'accroche de mesure du temps facultatifs~; \texttt{RetxServer}~; \texttt{recorded}.
Rust \texttt{firm\_feedhandler::Handler} avec les mêmes éléments.
\bfield{Règles} Chaque numéro de séquence est publié une seule fois, dans l'ordre~; les messages retenus ne sont libérés que dans l'ordre~; un trou attend l'autre
ligne, puis la retransmission (une demande par trou, réponses plafonnées par le serveur), puis un instantané~; le gestionnaire est
une fonction de ses entrées horodatées~; l'\omterm{def:ll:the-feed-handler:stale}{indicateur de carnet périmé} est levé exactement pendant qu'un trou est ouvert~; aucune allocation par paquet sans
trou.
\bfield{Tests d'acceptation} \texttt{code/firm/feedhandler/}~: sur le jeu de données commun, exécution propre, exécution avec retransmission et exécution avec instantané, avec
nombres d'événements, condensats, intervalles de péremption et compteurs identiques en Python, C++ et Rust~; l'exécution avec retransmission égale à
l'exécution propre, événement par événement~; l'exécution avec instantané égale l'exécution propre après l'instantané et dans son carnet final~; publication sur
\texttt{firm.ring} relue condensat par condensat~; arbitrage sur des paquets scénarisés.
\bfield{Pour aller plus loin} Une seule demande pour tous les trous d'un même trou~; reprise par instrument sur un flux multi-instruments~; lecture des lignes
depuis des sockets (\omterm{def:ll:linux-tuning:polling}{attente active} du chapitre~13) avec la marque de pic de la file de réception comme compteur.
\end{build}

\begin{omsources}
\item Page de manuel Linux \texttt{socket(7)} (\texttt{SO\_RCVBUF})~; Linux kernel documentation, ``Documentation for
/proc/sys/net/'' (\texttt{rmem\_max}, \texttt{rmem\_default}, \texttt{netdev\_max\_backlog}).
\item OPRA, \emph{Key Operating Metrics of U.S. Options Securities Information Processor} (août 2026)~; SIAC, \emph{Revised
OPRA Capacity Projections} (septembre 2025).
\item CME Group Client Systems Wiki, MDP~3.0: ``Market Data Snapshot -- Full Recovery''.
\item Les protocoles et services du simulateur, \path{code/firm/exchsim/PROTOCOL.md} (livre~10, chapitre~26).
\end{omsources}

\section{Exercices}

\begin{exercise}[$\star$]\label{exo:ll:the-feed-handler:1}
Le gestionnaire attend le numéro de séquence 1\,000 et reçoit des paquets commençant à 1\,003 (3 messages) sur la ligne A, puis 1\,000 (3
messages) sur la ligne B. Que publie-t-il, dans quel ordre, et quel est l'intervalle de péremption si le second paquet arrive
\qty{40}{\micro\second} après le premier~?
\end{exercise}

\begin{exercise}[$\star$]\label{exo:ll:the-feed-handler:2}
Un gestionnaire demande un tampon de réception de 256~KiB. Que renvoie \texttt{getsockopt}, et combien de datagrammes de 64~octets le
tampon peut-il contenir, selon la mesure de la \cref{fig:ll:the-feed-handler:rcvbuf}~?
\end{exercise}

\begin{exercise}[$\star$]\label{exo:ll:the-feed-handler:3}
Avec des instantanés toutes les 500 millisecondes et un instantané dont la réception prend 5 millisecondes, quelles sont les \omterm{def:ll:the-feed-handler:stale}{latences de reprise}
moyenne et au pire d'un trou non récupérable~?
\end{exercise}

\begin{exercise}[$\star\star$]\label{exo:ll:the-feed-handler:4}
Pourquoi le gestionnaire attend-il l'autre ligne avant de demander une retransmission~? Que coûte une attente plus courte, et une
attente plus longue~?
\end{exercise}

\begin{exercise}[$\star\star$]\label{exo:ll:the-feed-handler:5}
Une rafale arrive à 2 millions de datagrammes par seconde pendant 30 millisecondes~; le gestionnaire en traite 1.2 million par seconde. Combien de
datagrammes se mettent en file d'attente, et quel tampon de réception (demandé) les contient s'il est facturé 960 octets par datagramme~?
\end{exercise}

\begin{exercise}[$\star\star$]\label{exo:ll:the-feed-handler:6}
Après une reprise par instantané, le flux d'événements du gestionnaire diffère de celui de l'exécution propre. En quel sens la sortie est-elle néanmoins
«~identique une fois reprise~», et comment la construction le teste-t-elle~?
\end{exercise}

\begin{exercise}[$\star\star\star$]\label{exo:ll:the-feed-handler:7}
\emph{Programmation.} Faire demander au gestionnaire tous les trous d'un même trou en une seule demande (une liste d'intervalles), dans la référence Python et
le gestionnaire C++, et vérifier sur la minute de la \cref{fig:ll:the-feed-handler:minute} que la péremption maximale des retransmissions baisse.
\end{exercise}

\begin{exercise}[$\star\star\star$]\label{exo:ll:the-feed-handler:8}
\emph{Trouver la faille.} «~Quand notre gestionnaire détecte un trou, il abandonne tout ce qu'il détient et attend le prochain instantané~; ainsi
le carnet est toujours cohérent.~»
\end{exercise}

\section{Problème~: Le pic d'ouverture}

\begin{problem}[{Problème du week-end --- tampons pour l'ouverture, péremption quand ils cèdent}]\label{pb:ll:the-feed-handler:1}
Un gestionnaire lit une seule ligne d'un flux dont la rafale d'ouverture atteint 1.5 million de datagrammes par seconde pendant 20 millisecondes, un
message par datagramme~; le gestionnaire traite un datagramme en \qty{2}{\micro\second} (décodage, mise à jour du carnet, publication). Prendre la
charge du noyau par datagramme de 64~octets et son doublement de \texttt{SO\_RCVBUF} tirés de la \cref{fig:ll:the-feed-handler:rcvbuf}, un
\texttt{rmem\_max} maximal de 4~MiB, et un instantané chaque seconde.

\medskip\noindent\textbf{Partie I --- La file.}
\begin{enumerate}
\item Quel est le débit de service du gestionnaire, et quel est son taux d'utilisation pendant la rafale~?
\item Combien de datagrammes sont en file à la fin de la rafale~?
\item Combien de temps après la rafale la file met-elle à se vider si le débit tombe à 200\,000 par seconde~?
\item Quel est le délai de file d'attente du dernier datagramme de la rafale~?
\end{enumerate}

\medskip\noindent\textbf{Partie II --- Le tampon.}
\begin{enumerate}[resume]
\item Combien d'octets de tampon de réception le noyau doit-il accorder pour contenir la file~?
\item Que faut-il demander avec \texttt{SO\_RCVBUF}, et cela est-il permis sur cette machine~?
\item Avec le plus grand tampon autorisé, combien de datagrammes sont perdus~?
\item Qu'est-ce qui d'autre pourrait contenir la rafale~?
\end{enumerate}

\medskip\noindent\textbf{Partie III --- Quand cela échoue.}
\begin{enumerate}[resume]
\item Le serveur de retransmission peut-il combler la perte~?
\item Si non, quelle péremption en découle, en moyenne et au pire~?
\item Pendant cette péremption, combien de datagrammes arrivent et doivent être retenus~?
\item Que doit faire la stratégie dans l'intervalle~?
\end{enumerate}

\medskip\noindent\textbf{Partie IV --- Le verdict.}
\begin{enumerate}[resume]
\item Énoncer le \emph{résultat nommé}~: le tampon de réception qui absorbe la rafale sans perte, et la péremption qu'implique
un intervalle d'instantanés lorsqu'il y a perte.
\item De combien le gestionnaire devrait-il être plus rapide pour n'avoir besoin d'aucune file~?
\item Pourquoi lire les deux lignes double-t-il le problème~?
\item Que dit le pic à la milliseconde de l'OPRA d'un gestionnaire mono-fil pour un flux d'options complet~?
\item Que mesurerait-on en production pour dimensionner le tampon~?
\item Quel compteur indique qu'une rafale a approché la limite~?
\item Pourquoi l'intervalle d'instantanés est-il une décision commerciale de la plateforme, avec un coût de latence pour la société de gestion~?
\item En une phrase~: que fait un \omterm{def:ll:the-feed-handler:handler}{feed handler} lorsqu'il ne peut pas suivre~?
\end{enumerate}
\end{problem}

\section{Questions d'entretien}

\begin{interviewq}[$\star$ \iqroles{développeur}]\label{iq:ll:the-feed-handler:1}
Que fait un \omterm{def:ll:the-feed-handler:handler}{feed handler} de données de marché, de la carte réseau à la stratégie~?
\end{interviewq}

\begin{interviewq}[$\star\star$ \iqroles{développeur}]\label{iq:ll:the-feed-handler:2}
Un trou apparaît sur la ligne A. Décrivez ce que fait le gestionnaire jusqu'à ce que le carnet soit de nouveau correct.
\end{interviewq}

\begin{interviewq}[$\star\star$ \iqroles{développeur}]\label{iq:ll:the-feed-handler:3}
Comment dimensionneriez-vous le tampon de réception d'un socket UDP pour un flux~? Que fait le noyau de la valeur que vous fixez~?
\end{interviewq}

\begin{interviewq}[$\star\star$ \iqroles{développeur}]\label{iq:ll:the-feed-handler:4}
Comment savez-vous, en aval, que le carnet que vous lisez est périmé~?
\end{interviewq}

\begin{interviewq}[$\star\star$ \iqroles{développeur}]\label{iq:ll:the-feed-handler:5}
Pourquoi normaliser les messages de chaque plateforme en un seul enregistrement interne~? Que coûte cette normalisation~?
\end{interviewq}

\begin{interviewq}[$\star\star\star$ \iqroles{développeur, chercheur}]\label{iq:ll:the-feed-handler:6}
Comment testeriez-vous les chemins de reprise d'un \omterm{def:ll:the-feed-handler:handler}{feed handler} avant la production, et prouveriez-vous qu'il retrouve le bon carnet~?
\end{interviewq}
